\section{Founder-Market Fit: The Right Founder for Which Market?}

\subsection{A Task-Specific Definition}

Founder-market fit is often used as an aura surrounding a persuasive founder story. Someone has an impressive background, knows the industry's language, and seems emotionally committed. Those things can be meaningful, but they are not a validated formula for future business success.

I would define founder-market fit more narrowly:

\begin{tcolorbox}[
  enhanced,
  colback=BayesPale!25!white,
  colframe=BayesBorder!60!white,
  boxrule=0.4pt,
  arc=2pt,
  left=6pt,right=6pt,top=4pt,bottom=4pt,
  before skip=3pt,after skip=4pt
]
\footnotesize\color{BayesInk}
\textbf{Founder-market fit is the degree to which a founder's currently deployable capabilities, knowledge, access, resources, and motivation match the critical tasks of discovering and building a business in a particular market.}
\end{tcolorbox}

Notice the word \textit{critical}. A founder need not excel in everything. They must either possess or acquire the capabilities needed to clear the most important bottlenecks at the current stage.

A market can be attractive without being attractive \textbf{for you now}. And a founder can fit a market extremely well while the business opportunity within it is economically weak.

\subsection{5 Dimensions of Fit Plus Motivation}

\begin{itemize}[leftmargin=11pt,itemsep=1.5pt,topsep=1.5pt,parsep=0pt]
  \item \textbf{Problem fit:} Do you understand the problem deeply enough to distinguish a genuine failure from an imagined inconvenience?
  \item \textbf{Customer fit:} Do you understand the buyer's incentives, language, constraints, and procurement journey? Can you get real meetings?
  \item \textbf{Solution fit:} Can you build or orchestrate a solution that works, with appropriate quality and risk controls?
  \item \textbf{Distribution fit:} Can you earn attention and trust from intended customers at a viable acquisition cost?
  \item \textbf{Economic and operating fit:} Can you deliver, support, finance, and improve the model under your current capital and time limits?
  \item \textbf{Motivational fit (the 6th lens):} Can you sustain the work without your energy disappearing when the market behaves differently from your fantasy?
\end{itemize}

\subsection{The Bottleneck Principle}

\textbf{The bottleneck principle: a founder does not need maximum ability in every dimension, but a severe deficiency in 1 indispensable dimension can dominate everything else.}

Imagine a fictional founder:

\begin{tcolorbox}[
  enhanced,
  colback=white,
  colframe=BayesBorder!80!white,
  boxrule=0.4pt,
  arc=2pt,
  left=4pt,right=4pt,top=4pt,bottom=4pt,
  before skip=3pt,after skip=5pt
]
\scriptsize
\setlength{\tabcolsep}{3pt}
\renewcommand{\arraystretch}{1.18}
\begin{tabularx}{\linewidth}{@{} >{\bfseries\color{BayesInk}}p{3.2cm} >{\color{BayesSecondary}}X @{}}
\toprule
\textbf{Fit Dimension} & \textbf{Founder Judgment} \\
\midrule
Technical solution capability & Strong, demonstrated in production. \\
Understanding customer problem & Moderate, needs workflow interviews. \\
Buyer access & Weak, zero executive contacts. \\
Survival through procurement & Weak, cannot endure 12-month sales cycle. \\
Motivation for outbound sales & Weak, actively avoids cold calling. \\
\bottomrule
\end{tabularx}
\end{tcolorbox}

If this person wants to sell enterprise systems with an 18-month procurement cycle, technical superiority may not help enough. The bottleneck is access and survival. A smaller, faster-selling segment (or a distribution partner) might be a much better starting point.

\subsection{A Provisional 6-Dimension Scorecard}

Score each dimension as follows:
\begin{itemize}[leftmargin=11pt,itemsep=1pt,topsep=1pt,parsep=0pt]
  \item \textbf{0:} Absent or incompatible; no demonstrated capability or access.
  \item \textbf{1:} Plausible but untested; belief or adjacency.
  \item \textbf{2:} Demonstrated in a nearby context; transfer uncertain.
  \item \textbf{3:} Demonstrated in target context; verified repeated evidence.
\end{itemize}

\begin{tcolorbox}[
  enhanced,
  colback=white,
  colframe=BayesBorder!80!white,
  boxrule=0.4pt,
  arc=2pt,
  left=3pt,right=3pt,top=3pt,bottom=3pt,
  before skip=3pt,after skip=5pt
]
\scriptsize
\setlength{\tabcolsep}{2pt}
\renewcommand{\arraystretch}{1.18}
\begin{tabularx}{\linewidth}{@{} >{\bfseries\color{BayesInk}}p{2.2cm} *{6}{c} >{\color{BayesTeal}}p{1.7cm} @{}}
\toprule
\textbf{Opportunity} & \textbf{Pr} & \textbf{Cu} & \textbf{So} & \textbf{Di} & \textbf{Ec} & \textbf{Mo} & \textbf{Bottleneck} \\
\midrule
AI Eval Consulting & 3 & 2 & 3 & 2 & 2 & 3 & Repeatable lead gen \\
Skincare D2C Brand & 0 & 0 & 1 & 0 & 0 & 1 & Domain/supply \\
Tech Cohort for Devs & 2 & 2 & 3 & 1 & 2 & 3 & Acquisition/trust \\
\bottomrule
\end{tabularx}
\end{tcolorbox}

\textbf{Illustrative scores only.} Avoid comparing summed scores mechanically; the weakest mandatory dimension matters far more than the sum. Label any dimension requiring legal authorization, unavailable weekday capacity, or unaffordable capital as a \textbf{gate}, not a low score that can be averaged away.
